\section{Wednesday 2/4: Welcome, what is statistical learning, basics of supervised learning}

\subsection{What is statistical learning?}
\label{sec_veryintro}

Statistical learning very broadly refers to a set of tools for \emph{learning} something from data. 

Consider the following three scenarios:
\begin{itemize}
\item \textbf{Scenario 1:} A researcher comes up with the hypothesis that college students who sleep more than 7 hours per night have higher GPAs, on average, than those who sleep fewer than 7 hours per night. She goes out and collects a random sample of students, and asks each of them (1) if they sleep more than 7 hours per night, and (2) for their GPA.

She then fits the linear model

$$
\widehat{GPA} = \hat{\beta}_0 + \hat{\beta}_{sleep} (\mathrm{sleep\_more\_than\_7}),
$$
and uses the estimated value $\hat{\beta}_{sleep}$, along with its associated p-value and confidence interval in R, to answer her research question.
\item \textbf{Scenario 2:} A researcher is interested in what factors impact the GPA of a college student, so she collects a random sample of students and writes down a lot of information about them. She searches through the data, tries different models, looks at plots, and ends up realizing that “hours of sleep per night” is an important factor in determining GPA. As the relationship between hours of sleep per night and GPA does not look linear, she makes a new binary variable: yes/no do you sleep more than 7 hours per night.

At the end of this process, she fits the linear model
$$
\widehat{GPA} = \hat{\beta}_0 + \hat{\beta}_{sleep} (\mathrm{sleep\_more\_than\_7}).
$$
\item \textbf{Scenario 3: } A researcher is interested in predicting the GPA of a college student based on various factors. She collects a random sample of students and writes down a lot of information about them. Without bothering to look at plots, she asks R to fit her a very complicated model that takes in all possible input variables $X_1,\ldots,X_p$ and allows them to contribute non-linearly and in interacting ways to produce an estimate $\widehat{GPA}$. She sees that, on her dataset, her predictions are quite accurate ($R^2=0.98$). 
\end{itemize}

Based on our definition, all three of these scenarios are examples of statistical learning. In scenario 1, we \emph{learned} the value of $\hat{\beta}_{sleep}$: this was an estimate of a pre-specified parameter $\beta_{sleep}$. In scenario 2, we \emph{learned} not only $\hat{\beta}_{sleep}$, but we also learned to be interested in the parameter $\beta_{sleep}$ via a process of \emph{variable selection} (choosing ``sleep") and \emph{feature engineering} (deciding to make it binary). Finally, in scenario 3, we \emph{learned} a prediction function, which could be applied to new datapoints. 

While all three of these scenarios are examples of statistical learning, only scenario 1 is an example of classical statistics. Classical statistics  assumes that our data analysis plan is completely pre-specified before looking at our data; this is the only setting in which classical confidence intervals and p-values are formally valid. When we talk about statistical learning, we are often focusing on methods and strategies that move beyond classical statistics: we are often discovering patterns in complex data without a specific, pre-specified plan. Statistical learning is often much more exploratory than classical statistics!\\
\\
There is nothing ``right" or ``wrong" about the approaches in scenarios 1-3; it all depends on the goals of the researcher. In scenario 1, we get to rigorously test a hypothesis using tools we love (p-values!), but we sacrifice the flexibility to explore our data and uncover new ideas. In scenario 2, we get to be a lot more flexible and adaptive, but we give up our rigorous hypothesis testing, and manually trying so many different models might get tedious. If our goal is truly prediction, rather than inference, perhaps we like scenario 3, which lets us build a maximally flexible model without too much manual effort. \\
\\
With this very broad definition of statistical learning in mind, we can now get into some specifics.

\subsection{Basic terminology, and types of learning}

In statistical learning, we use $n$ to denote the number of observations. For $i=1,\ldots,n$, we measure $p$ \emph{covariates} or \emph{predictors}, $x_{i1}, x_{i2}, \ldots, x_{ip}$. We denote the entire matrix of observed covariate with $\bold{X}$, which has dimensions $n \times p$. I will typically use capital, bolded letters to denote matrices.

In supervised learning, we also measure a response variable $y_i$ for each observation. I will typically use bold lowercase letters to denote vectors of length $n$, so we could collect these observations all together to make $\bold{y}$.

Notationally, I will use capital, unbold letters to denote random variables, which could be either vectors or scalars. I use lowercase letters for observed values. In our example from Section~\ref{sec_veryintro}, we let $X_1$ denote the number of hours that a randomly selected student sleeps per night and let $Y$ denote their GPA. Once we observe a particular $n$ students for our dataset, we let their sleep values be $\bold{x}_1 = (x_{11}, x_{12},\ldots, x_{1n})^\top$ and we let their GPA values be $\bold{y} = (y_1,\ldots,y_{n})^\top$.  

Putting it all together, the data that we start with for a supervised learning problem is a matrix $\bold{X} \in \mathbb{R}^{n \times p}$ and a response vector $\bold{y} \in \mathbb{R}^n$, and we are generally interested in using the predictor variables to explain or predict the response variable. We usually hope that whatever we learn would generalize to hypothetical new realizations $(X,Y)$. 

In contrast, the setting with no response variable $y$ where we just want to uncover patterns in $\bold{X}$ is known as \emph{unsupervised learning}.

In this course, we will spend around 9 weeks on supervised learning, and only 2 weeks at the end on unsupervised learning. However, I do not want to give you the impression that unsupervised learning is less important! Unsupervised learning is extremely important, and is in fact the back-bone of recent generative AI technology. We start with supervised learning because it has a more natural connection to classical statistics and builds more naturally on your previous work. Hopefully, some of the concepts that you learn during our supervised learning unit will help you be equipped to delve more deeply into unsupervised learning in the future if you choose.

 There is one more type of learning that has become really important in recent years. This is the field of \emph{semi-supervised learning}, in which we have a (potentially small) set of \emph{labeled data} (data with predictors + response) and a potentially much larger set of \emph{unlabeled data} (data that is missing the response). While we will not cover this in lecture, there are some great potential final project topics related to semi-supervised learning. 
  
\subsubsection{A side-bar about counting n and p}

We didn't talk about this in class, but I want to quickly make the point that counting $n$ and $p$ is not always trivial. 

Consider a dataset with one categorical variable that takes on values $\{ \text{cat}, \text{dog}, \text{hamster}, \text{rabbit} \}$. I could represent this variable using an $\bold{X}$  with one column that stores $\{$dog, dog, cat, dog, hamster, cat, dog, $\ldots \}$. However, for the purpose of model fitting, it is likely going to be convenient to instead let $\bold{X}$ be a matrix with four columns, that store binary variables indicating ``is this a dog", ``is this a cat", etc. For the purposes of parameter-counting, this is now a dataset with $p=4$, not $p=1$. 

Similarly, suppose we take measurements on $920$ students, who are spread out between $8$ schools. The schools are randomly assigned to either have required yoga PE class, or to have traditional PE class, and we measure resulting test scores at the end of the year. While we have $920$ students, the students within a given school are not \emph{independent} of one another. On the other hand, the schools underwent random assignment and can be safely assumed to be independent of one another. Thus, is the sample size $n$ the $920$ students or the $8$ schools? It turns out that for this \emph{clustered} or \emph{hierarchical} data, we have a notion of \emph{effective sample size} that is somewhere between $920$ and $8$, and depends on how correlated the students are within the schools. 

We should also mention that data does not always come to you in a form where it looks like $\bold{X} \in \mathbb{R}^{n \times p}$ and $\bold{y} \in \mathbb{R}^n$. Consider the task of image classification. You get a set of $500$ images, which are stored as $400 \times 400$ pixel images, where each pixel records a numerical value between 0 and 255 for red, green, and blue. Each image is associated with a label in $\{\text{cat},  \text{dog}\}$, and you observe $y_1,\ldots, y_{500}$. In this case, the RGB values in the pixels in the images are your predictor values, but they don't come in a simple $n \times p$ matrix. We need to reshape the data so that each image $i$ has an associated vector $x_i \in \mathbb{R}^{400 \times 400 \times 3}$ that stores all three color values for all $400 \times 400$ pixels. So $p$ is very large! It will also likely be convenient to code $y_i$ such that a $1$ denotes a dog and a $0$ denotes a cat, such that $y$ becomes numerical. When there are more than $2$ classes, such as $\{\text{cat},  \text{dog}, \text{fish}\}$, we will need more than one column to store our response variables numerically. 

Who knew that counting $n$ and $p$ could be so complicated! In this class, just think of it as the number of rows and columns (minus the response columns) of the data, and remember these subtleties in case they ever come up. \\

\subsection{Goals of supervised learning}

Generally in supervised learning, we assume that there is some relationship between $Y$ and $X$, such that we can write $Y$ as
$$
Y = f(X)+\epsilon.
$$
The function $f()$ is a fixed but unknown function that represents the systematic information that $X$ provides about $Y$, and $\epsilon$ is a random error term, which should satisfy the condition that $E[\epsilon] = 0$ and that $\epsilon$ is independent of $X$. These conditions ensure that $f()$ is capturing all parts of $Y$ that can be explained by $X$, and $\epsilon$ captures the parts of $Y$ that cannot be explained by $X$. 

The main task of supervised learning is: given $(x_i,y_i)$ for $i = 1,\ldots,n$, can we estimate $f$? Meaning, can we come up with a function $\hat{f}()$ such that, $Y \approx \hat{f}(X)$? 

When $Y$ is numerical, this task is called regression. When $Y$ is categorical, this task is called classification. 

There are two reasons why we might want to come up with a function $\hat{f}$.
\begin{enumerate}
\item \textbf{Prediction:} We might want to be able to come up with predictions 	$\hat{Y} = \hat{f}(X)$ for future realizations of data where $Y$ itself is not observed.
\item \textbf{Interpretation:} We might want to understand which predictors in $X$ are most associated with $Y$ in order to draw a scientific conclusion, or at least take a step towards drawing a scientific conclusion. 
\end{enumerate}
As we will see throughout the semester, these two goals can sometimes be at odds with one another! If we only care about prediction, we do not need our function $\hat{f}()$ to be interpretable: we can use an extremely complex model (known as a ``black box"), and we will be happy as long as $\hat{Y} \approx Y$. On the other hand, if we want to draw scientific conclusions, then we might wish to use a simple, interpretable model for $\hat{f}()$ that lends itself to rigorous inference. We should also note that it is not always an either-or situation. Sometimes, we want to make good predictions and also understand which variables are contributing significantly to the predictions. Managing the tradeoff, or lack thereof, between these two goals will be one of our fundamental themes for the semester. When starting a supervised learning problem, always think carefully about which of these goals (or both) is important to you! 

\subsection{What methods do you already know, and why will we move beyond them?}

You all already know at least one method for regression (linear regression) and one method for classification (confusingly, named logistic regression). While these methods are taught in the most classical of statistics courses, they are certainly examples of machine learning methods! 

Why will this class move beyond these methods? 

Let's focus on the linear regression case, and suppose that we have $p$ predictors. This means that we will let:
$$
\hat{f}(X) = \hat{\beta}_0 + \hat{\beta}_1 X_1 + \ldots + \hat{\beta}_p X_p
$$

Over the next few weeks, we will focus on ways to make this model both \emph{more} and \emph{less} flexible. You have actually seen both of these concepts in Stat 346. 

\begin{itemize}
\item We need ways to make this more flexible, because this simple linear form might not do a good job approximating $Y$. What if the relationship between $X_1$ and $Y$ is quadratic, or is otherwise wiggly? What if we need interaction terms between variables? 
\item We also need ways to make this less flexible. You all know from Stat 346 that, if $p$ is very large, we run into trouble. If $p > n$, there is no unique least squares solution to the linear regression problem. If $p < n$ but is otherwise large, fitting this linear regression is probably a bad idea: the model is so complex that we risk overfitting to our training set. If our predictors are redundant, the model can be very unstable. 
\end{itemize}

You learned in Stat 346 that you can always add transformations of variables, or interaction terms, to make a model more flexible and still fit it in linear regression. You also learned that variable selection techniques or PCA can help reduce the complexity of a model to avoid overfitting or instability. 

We will start by talking about these concepts (more and less flexible models) mainly in the context of things that at least look like linear regression. However, there is also no rule that says $\hat{f}()$ has to be a linear function! 

In recent decades, as datasets have become larger and computers have become more powerful, there have been a proliferation of complex methods for estimating $\hat{f}()$. Some of these look like fancy linear regression, but some look entirely different! Why? 

\begin{itemize}
\item Datasets have gotten bigger. As we will learn in this class, complex models are often only appropriate when applied to enough data. Thus, it is natural that the ``big data" era has gone hand-in-hand with the development of complex statistical learning methods. 
\item Computers have gotten better, faster, and cheaper, which has enabled the fitting of complex models. 
\item Free, open-source software programs like R and python have allowed researchers to develop algorithms and then easily share them with non-experts. This has allowed many methods to become popular. 
\end{itemize}

Given the huge number of statistical learning methods, we could spend each class this semester learning a new method, such that at the end you are left with a huge cookbook of methods to choose from for your future data analysis needs. In fact, this class will be a little bit like this. But, there is a problem with this approach of treating statistical learning like a cookbook. New machine learning methods get published every day, and the state-of-the-art is constantly changing. If your goal is to learn a list of methods or algorithms, you will find yourself continually behind and out-of-date! Thus, the goal of this class is not really to teach you a list of methods. Instead, the goal is to teach you the fundamental concepts and overarching themes that go into the development of the methods, such that you can compare methods, recognize their limitations, and learn new methods on your own in the future. This goal helps explain two features of this course: 
\begin{itemize}
	\item We will spend a fair amount of time on older, less state-of-the-art methods. We will use these methods to illustrate general principles and themes, even if they may not be the methods that you will use in practice in your future. 
	\item You will all be doing an \emph{algorithm presentation} this semester. You will be responsible for doing independent reading to learn about a method. You will need to synthesize what you learn, and give a 20 minute presentation to the class explaining the method. This skill mirrors the way you will all interact with machine learning in the future: you will be continually learning and synthesizing new methods. 
\end{itemize}



